\section{Introduction}\label{sec:intro}

A word program may inspect a count, use a shift, propagate a carry, and finally return only a few facts about its result. A proof that reconstructs every intermediate value can retain substantially more information than the consumer needs. QKF asks which observations determine this particular action, and which of their values are forced by the supplied information and the operations already justified on the carrier.

This is Paper III of a series by Leonid Shcherbakov. Paper I develops initial semantics for incomplete actions, forced domains, kernel saturation, and completion~\cite{ShcherbakovI}. Paper II studies equality visibility through finite ground interfaces and distinguishes congruence-proof depth from the size of sequential rewriting contexts~\cite{ShcherbakovII}. The present paper connects these distinctions to executable proof objects. Native semantics supplies the facts to which the algebra applies; a quotient is useful only when its labels and its interaction with the consumer have been justified.

The first version described an exact finite algebraic kernel and a KnownBits certifier whose final language was coordinatewise: nine one-bit rows suffice for a theorem at every positive width. Version 2 added contextual observations, labeled whole-word reconstruction, signed mask--interval carriers, and a Lean pilot. Version 3 develops two complementary paths. One derives a sufficient finite source interface and exposes its checked action. The other obtains native facts for an independently specified consumer and proves the resulting closed equality directly. Their common output is a scoped, replayable guarantee; their discovery and execution costs need not be the same.

\paragraph{Contributions and stable claims.}
Claims K1--K10 retain their meanings: finite ground closure and its sparse contract (K1--K3), coordinatewise source replay and semantic classification (K4--K5), sufficient labeled observations and finite gluing (K6--K7), the upper and lower masked-bound contracts (K8--K9), and the original Lean pilot (K10). This revision adds executable initial pure-row semantics, restricted automatic source-interface inference, chronological typed ground proofs, checked-lemma composition, and a summary SDK. The new formal results concern the decoded positive proof calculus and its composition; they do not enlarge K10 into a proof of the entire software stack.

\begin{center}\small
\begin{tabularx}{\textwidth}{@{}l>{\raggedright\arraybackslash}X>{\raggedright\arraybackslash}X@{}}
\toprule
Layer & Computation & Evidence in this revision\\
\midrule
Finite algebra & Forced relations, initial rows, equality visibility & Mathematical arguments, proof/model replay, independent finite checks.\\
Coordinatewise CLI & Source normalization and nine-row targets & Stable package contract and retained regressions.\\
Source interfaces & Observation inference, coverage, row execution & Source-bound packages and independent target products.\\
Consumer summaries & Native facts, ground DAGs, checked lemmas & Scoped SDK objects, portable proofs, controlled route comparisons.\\
Lean projects & Word fragments, typed ground soundness, lemma composition & Kernel-checked proofs with explicit assumptions.\\
\bottomrule
\end{tabularx}
\end{center}

The installed command-line package does not automatically acquire every research frontend. Paper II does not confer a minimum proof-depth guarantee on an ordinary normalization trace. Likewise, a successful factorization is not a consumer proof, and a finite separating model of insufficient premises is not necessarily a counterexample to the source program. The explicit maps between these layers are part of the specification.

\paragraph{Version and source identity.}
This is manuscript version 3.0, dated 22 September 2026, accompanying the \texttt{0.3.0a1} research-alpha release preparation. The technical and experimental snapshot is the merged main revision through PR51:
\begin{center}\small\texttt{eca90f9252da4480d77c1f8cb4b7062470c57b89}.\end{center}
The repository is \url{https://github.com/Len989/QKF-Certifier}. Artifact links in the new sections resolve to this immutable revision; individual reports additionally identify the engine that was measured. Earlier coordinatewise results retain their original baseline \texttt{f93561682319e8b911ed32c01583f2a496dc59fc}. The release manifest identifies the packaged revision and files. Manuscript version, software version, certificate schema, and semantic contract are separate identifiers. The prior Paper III deposit is version 2.0~\cite{QKFPaperIIIv2}; no new version DOI is presumed here.

The immediate research objective is testable: infer adequate interfaces within explicit source contracts, obtain only the facts a consumer needs, preserve their dependencies across reuse, and reduce total delivery and checking cost. The broader objective of handling substantially wider verification tasks remains open. The internal ordinary-SDK route is part of QKF: its favorable results inform implementation choices within this architecture.
